\documentclass[10pt]{article}
\usepackage[a4paper,margin=1.7cm]{geometry}
\usepackage[T1]{fontenc}
\usepackage{lmodern}
\usepackage{amsmath,amssymb}
\usepackage{longtable,booktabs,array}
\usepackage{listings}
\usepackage[hidelinks]{hyperref}
\lstset{basicstyle=\ttfamily\footnotesize,breaklines=true,columns=fullflexible,keepspaces=true,upquote=true}
\setcounter{MaxMatrixCols}{20}
\allowdisplaybreaks
\setlength{\parindent}{0pt}\setlength{\parskip}{4pt}\setlength{\emergencystretch}{3em}\sloppy\raggedright
\begin{document}

{\raggedright\LARGE\bfseries Symbols used in the README of Pre-gravity6plane\par}\medskip

{\small\begin{longtable}{>{\raggedright\arraybackslash}p{\dimexpr 0.4800\linewidth-2\tabcolsep\relax}>{\raggedright\arraybackslash}p{\dimexpr 0.4800\linewidth-2\tabcolsep\relax}}\toprule
symbol & definition \\\midrule
$x_0..x_7$ & coordinates: hidden, space (1-3), time (4), superluminal time (5-7); the Pre-Universe part writes them x0, ..., x7 \\
$\eta$ & $\mathrm{diag}(1,1,1,1,-1,-1,-1,-1)$ \\
$\tau_A$, $\bar\tau_A$ & real 8x8 matrices of Pre-gravity6plane.nb, $\bar\tau_A = \sigma\tau_A^T\sigma$, $\sigma = \begin{pmatrix}0&I_4\\ I_4&0\end{pmatrix}$ \\
$T^A = (T16^A)[A]$ & $\begin{pmatrix}0&\bar\tau_A\\ \tau_A&0\end{pmatrix}$, real Dirac matrices \\
$S^{AB}$ & $\frac14[T^A,T^B]$ \\
Spin(4,4), Pin(4,4) & the spin group of $\eta$, whose Lie algebra is spanned by the $S^{AB}$; the pin group, generated by Spin(4,4) and the reflections $T^C$ \\
$\Gamma$, $P_{L,R}$, $\sigma_{16}$ & $T^0\cdots T^7$; $\frac12(I_{16}\mp\Gamma)$; $T^0T^1T^2T^3$ \\
$\psi$, $\bar\psi$ & dirac16complex (complex Grassmann 16-spinor) $\psi = (\psi_0,\dots,\psi_{15})$, $\bar\psi = \psi^\dagger\sigma_{16}$; in the Pre-Universe part, a real 16-component spinor $\psi\in\mathbb R^{16}$ \\
$\psi^\dagger$, $\psi^*$ & conjugate transpose and complex conjugate of $\psi$ \\
$\psi_1$, $\psi_2$ & type-1 and type-2 halves of $\psi$: components $0..7$ and $8..15$; in the Pre-Universe part, $\psi_1 = \psi_1^a$ and $\psi_2 = \psi_2^{(a)}$ ($a = 0, \dots, 7$), the $\overline{SO(4,4;\mathbb R)}$ counterparts of the left-handed and right-handed spin(8) spinors, with $\psi\in\mathbb R^{16} = \{\psi_1\in\mathbb R^8\oplus\psi_2\in\mathbb R^8\}$ \\
$\partial_A$, $\partial_\mu$ & $\partial/\partial x_A$, $\partial/\partial x_\mu$ \\
$\mathcal L$, $m$, $\lambda$ & Lagrangian, mass, quartic coupling; $\mathcal L_m$ and $\mathcal L_{-m}$ are $\mathcal L$ with mass parameter $m$ and $-m$ \\
$K$, $K_{ab}$, $\delta$ & $K = -i\sigma_{16}T^4$ with entries $K_{ab}$; $\{\psi_a(x),\psi_b^\dagger(y)\} = K_{ab}\delta(x-y)$ at equal $x_4$, with $\delta$ the Dirac delta of the other seven coordinates \\
$\psi_\pm$, $S_\pm$, $\mathcal L_D$ & $K\psi_\pm = \pm\psi_\pm$; $S_\pm = \psi_\pm^\dagger(-iT^4)\psi_\pm$; $\mathcal L_D[\chi] = \chi^\dagger(-iT^4)(T^A\partial_A - m)\chi$ \\
$M$, $H$ & mass and Hubble parameters of Pre-gravity6plane.nb; $m = -2M$ \\
$\Psi_{16}$ & Universe wave function of Pre-gravity6plane.nb, $\Psi_{16} = (f16[0][x_0,x_4], \dots, f16[15][x_0,x_4])$; the Pre-Universe part writes it [CapitalPsi]16, the 16-component $\overline{O(4,4)}$ spinor of the Dirac equation for the Pre-Universe \\
$z$, $t$, $p$, $q$, $u$ & $z = 6Hx_0$, $t = Hx_4$, mode $e^{pz+qt}u$ with $u$ a constant spinor \\
$F_A{}^a$, $F_a{}^A$, $V_1$, $V_2$ & triality frames and vector forms: $V_1 = \eta(\psi_1^T\sigma F_A{}^a)$, $V_2 = F_a{}^A\psi_2$ \\
$n_k$, $Z$, $M^6$ & $n_k = V_k[\mathrm{Re}\,\psi]$; $Z\in\mathbb R^8$; $M^6 = \{Z : n_1\cdot Z = n_2\cdot Z = 0\}$ \\
$E$, $\binom82$ & energy of a universe for translations of $x_4$ (a pair of universes has energies $\pm E$); binomial coefficient, $\binom82 = 28$ \\
Re, Im & real and imaginary parts; for $x$ in Os, Re $x = \langle x,1\rangle$ and Im $x = x -$ Re $x$ \\
$f$, $r_s$, $H_\Lambda$ & $1-r_s/x_0$ (modern), $1-r_s/x_0-H_\Lambda^2x_0^2$ (future) \\
$a_4$ & function in the vielbein of Pre-gravity6plane.nb; $a_4(Hx_4) = Hx_4$ for the primordial field \\
$e^A{}_\mu$, $g_{\mu\nu}$, $g$, $ds^2$ & vielbein; metric $g_{\mu\nu} = \eta_{AB}e^A{}_\mu e^B{}_\nu$; $g$ its determinant, which vanishes on an Einstein-Rosen bridge; line element $ds^2 = g_{\mu\nu}dx_\mu dx_\nu$ (summed over $\mu, \nu = 0, \dots, 7$) \\
$E_A{}^\mu$, $\delta^\mu_\nu$, $\Gamma^\nu{}_{\mu\lambda}$, $\omega_\mu{}^A{}_B$ & inverse vielbein, $E_A{}^\mu e^A{}_\nu = \delta^\mu_\nu$ (Kronecker symbol); Christoffel symbols of $g_{\mu\nu}$; spin connection $\omega_\mu{}^A{}_B = e^A{}_\nu(\partial_\mu E_B{}^\nu + \Gamma^\nu{}_{\mu\lambda}E_B{}^\lambda)$ \\
$\gamma^\mu$, $D_\mu$ & $E_A{}^\mu T^A$; $\partial_\mu + \frac14\omega_\mu{}^A{}_BT_AT^B$ \\
$\rho_s$, $h_{KS}$, $\alpha$, $\beta$, $p$, $B$ & Kohn-Sham density, $\alpha p + BM\beta$, $-T^4T^0$, $-iT^4$, $-i\,d/d\rho$, lapse $e_4$ \\
$\Delta$SCF & excited states found as self-consistent Kohn-Sham solutions with a changed occupation of the levels \\
$\Omega(T)$ & Kohn-Sham grand potential at temperature $T$ \\
$R_1$ & $x_1\mapsto-x_1$ \\
$R$, $\kappa$, $T_H$, $\epsilon$ & scalar curvature, surface gravity, $\kappa/2\pi$, Killing energy \\
Kretschmann scalar & $R_{\mu\nu\rho\sigma}R^{\mu\nu\rho\sigma}$, with $R_{\mu\nu\rho\sigma}$ the Riemann tensor \\
$I_n$ & $n\times n$ identity matrix \\
$g_{jj}$, $b_{jj}$, $a$ & deflating metric functions $g_{jj} = b_{jj}e^{-ax_4}$ ($j = 5, 6, 7$), $a > 0$, $b_{jj} > 0$ \\
$w$, $y_j$ & coordinates of the deflating block: $w = \frac2ae^{ax_4/2}$, $y_j = \sqrt{b_{jj}}\,x_j$ \\
$\mathbb R^4$ & flat 4-space of $x_0, x_1, x_2, x_3$ \\
$\ell$, $p_5$, $x_4^*$ & affine parameter; $p_5 = -g_{55}\,dx_5/d\ell$; turning point $\frac1a\ln(b_{55}/p_5^2)$ \\
$\omega$, $k_0..k_3$, $q_5, q_6, q_7$ & frequency and wave numbers of a universe $e^{i(k_0x_0 + \dots + q_7x_7 - \omega x_4)}u$ (the indexed $q_j$ are wave numbers, not the exponent $q$); $u$ a constant spinor \\
$a_c$ & critical deflation rate $2H$ \\
$\chi$, $M$, $c_A$ & mode spinor with $\chi' = d\chi/dx_4 = M\chi$; $M = T^4(i\sum_{A\ne4}c_AT^A - m)$; real coefficients $c_A$ \\
$P_\pm$ & $\chi^\dagger\frac12(I_{16} \pm K)\chi$, weights of the +mass and -mass sectors \\
$\mathbb R^{4,4}$ & real eight-dimensional pseudo-Euclidean vector space with the $SO(4,4;\mathbb R)$-invariant quadratic form $x_0^2 + x_1^2 + x_2^2 + x_3^2 - x_4^2 - x_5^2 - x_6^2 - x_7^2$ \\
$SO(4,4;\mathbb R)$, $SO(4,4)$ & real linear maps of $\mathbb R^{4,4}$ with determinant 1 that preserve this quadratic form \\
$\overline{SO(4,4;\mathbb R)}$, $\overline{SO(4,4)}$ & universal cover of $SO(4,4;\mathbb R)$; the Pre-Universe part writes Spin(4,4) ( = $\overline{SO(4,4)}$ ) \\
Spin(3,4) & the subgroup of Spin(4,4) that preserves both $\mathbb R$ and Im Os; under it $\Psi_1$ and $\Psi_2$ transform the same way, and the index $(a)$ is written $a$ \\
SO(8), Spin(8), spin(8), $D_4$ & rotation group of $\mathbb R^8$ and its double cover; spin(8) spinors are the spinors of Spin(8); $D_4$ is their Dynkin diagram: a central node (adjoint representation) joined to the vector node and to the left-handed (type 1) and right-handed (type 2) spinor nodes \\
$O_{4,4}$, Os & split octonion algebra: real eight-dimensional algebra with unit 1, a non-degenerate indefinite inner product $\langle\ ,\ \rangle: Os\times Os\to\mathbb R$ (written \textless{} , \textgreater{} in the README), and a multiplication $Os\times Os\to Os$ with $\langle xy, xy\rangle = \langle x,x\rangle\langle y,y\rangle$ and $\langle 1,1\rangle = 1$; $\mathbb R$ is identified with the subalgebra generated by 1 \\
$Im\ Os$ & orthogonal complement of 1 in Os; $Os = \mathbb R\oplus Im\ Os$ \\
$x^*$ (x in Os) & conjugate, $x^* = 2\langle x,1\rangle - x$; x lies in Im Os if and only if $x^* = -x$ \\
$C_{4,4}$, id & (pseudo-) Clifford algebra generated by the $e^{4,4}_A$; id is its multiplicative identity \\
$e^{4,4}_A$ & the eight generators of $C_{4,4}$, $A = 0, \dots, 7$; they anticommute and satisfy $e^{4,4}_Ae^{4,4}_B + e^{4,4}_Be^{4,4}_A = id\ g_{AB}$, and the Pre-Universe part states that each has square $\pm$ id \\
$g_{AB}$, $g_{A,B}$ & $\mathrm{diag}(1,1,1,1,-1,-1,-1,-1)$ \\
$e_A$, $x^A$ & Minkowski vector basis of Os, $A = 0, \dots, 7$, with $e_0 = 1$ and $\langle e_A, e_B\rangle = g_{A,B}$; $x = x^Ae_A$ with real $x^A$ \\
$m^C_{AB}$ & multiplication constants, $e_Ae_B = e_Cm^C_{AB}$, constructed from the irreducible representation of the $C_{4,4}$ generators and a type 1 spinor u with $\langle u,u\rangle = 1$; they give the split octonion multiplication table (row $A$, column $B$: $e_Ae_B$) \\
$e^{1/2}_a$, $e^{1/2}_{(a)}$ & type 1 and type 2 spinor bases of Os, $a = 0, \dots, 7$ \\
$\Psi_1$, $\Psi_2$ & type 1 and type 2 spinors in Os: $\Psi_1 = \psi_1^ae^{1/2}_a$ and $\Psi_2 = \psi_2^{(a)}e^{1/2}_{(a)}$, with real $\psi_1^a$ and $\psi_2^{(a)}$ \\
$\mathcal C_{2n}$ & written Subscript[[ScriptCapitalC], 2 n]: complex Clifford algebra with $2n$ generators; Lord's reduced Brauer-Weyl generators form an irreducible representation of $\mathcal C_{2n}$ built from one of $\mathcal C_{2(n-1)}$ \\
tau8, tau16 & the real 8x8 and 16x16 matrices of Pre-gravity6plane.nb, obtained from Dirac's 4x4 gamma matrices with Lord's reduced Brauer-Weyl generators \\
r, 0, {\o}, t, u, m, g & in the passage quoted from Einstein and Rosen (1935): the coordinates r, $\theta$, $\phi$, t (printed r, 0, {\o}, t), the coordinate u with $u^2 = r - 2m$, the mass parameter m, and g, the determinant of the metric \\
$M_{4,4}$ & Minkowski spacetime of signature \{4,4\} (in the notation that the README asks the reader to ignore for now) \\
o, eA, ea, ea[Psi]1, ea[Psi]2, [Sigma], [Eta], X, EA & the notation that the README asks the reader to ignore for now, in Mathematica form: o, an element of $O_{4,4}$, expanded in the vector basis eA (components $o^j$) and, by Cartan's triality, in the type 1 and type 2 spinor bases ea[Psi]1 and ea[Psi]2; after the restriction to Spin(3,4) both spinor bases are written ea; \textless{}Subscript[ea, a], Subscript[ea, b]\textgreater{} = Subscript[[Sigma], ab] and \textless{}Subscript[eA, A], Subscript[eA, B]\textgreater{} = Subscript[[Eta], AB]; \{X, EA\} lies in $M_{4,4}\otimes O_{4,4}$ \\
\bottomrule\end{longtable}}

\end{document}